% \begin{table}[t]
% \renewcommand{\arraystretch}{0.95}
% \footnotesize
% \caption{Overall portfolio-level results on Dow-30 (25 assets, 2025 test). Mean $\pm$ std over 3 seeds.}
% \setlength{\abovecaptionskip}{0.1cm}
% \centering
% \setlength{\tabcolsep}{3pt}
% \begin{tabular}{lcccccc}
% \toprule
% Method & CR\%$\uparrow$ & SR$\uparrow$ & MDD\%$\downarrow$ & Calmar$\uparrow$ & Sortino$\uparrow$ & WR\%$\uparrow$ \\
% \midrule
% Buy \& Hold & ${\sim}$12.8 & ${\sim}$0.62 & ${\sim}$14.2 & ${\sim}$0.90 & ${\sim}$0.78 & ${\sim}$52.1 \\
% Equal-Weight & ${\sim}$10.5 & ${\sim}$0.55 & ${\sim}$12.8 & ${\sim}$0.82 & ${\sim}$0.69 & ${\sim}$51.4 \\
% SMA Cross & ${\sim}$8.2 & ${\sim}$0.48 & ${\sim}$10.5 & ${\sim}$0.78 & ${\sim}$0.61 & ${\sim}$53.8 \\
% \midrule
% PatchTST & ${\sim}$14.1 & ${\sim}$0.82 & ${\sim}$9.8 & ${\sim}$1.44 & ${\sim}$1.05 & ${\sim}$54.2 \\
% TimesNet & ${\sim}$13.5 & ${\sim}$0.78 & ${\sim}$10.2 & ${\sim}$1.32 & ${\sim}$0.98 & ${\sim}$53.6 \\
% iTransformer & ${\sim}$15.8 & ${\sim}$0.91 & ${\sim}$8.9 & ${\sim}$1.78 & ${\sim}$1.18 & ${\sim}$55.1 \\
% \midrule
% TradingAgents & ${\sim}$18.2 & ${\sim}$1.05 & ${\sim}$8.5 & ${\sim}$2.14 & ${\sim}$1.42 & ${\sim}$56.8 \\
% FinCon & ${\sim}$21.5 & ${\sim}$1.18 & ${\sim}$7.8 & ${\sim}$2.76 & ${\sim}$1.58 & ${\sim}$57.5 \\
% FinAgent & ${\sim}$16.4 & ${\sim}$0.95 & ${\sim}$9.2 & ${\sim}$1.78 & ${\sim}$1.25 & ${\sim}$55.4 \\
% R\&D-Agent & ${\sim}$19.8 & ${\sim}$1.08 & ${\sim}$8.1 & ${\sim}$2.44 & ${\sim}$1.45 & ${\sim}$56.2 \\
% AlphaAgent & ${\sim}$15.2 & ${\sim}$0.95 & ${\sim}$9.5 & ${\sim}$1.60 & ${\sim}$1.22 & ${\sim}$54.8 \\
% \midrule
% \textsc{FinOPD} (Ours) & $-$0.5 & $-$0.26 & 4.5 & $-$0.06 & $-$0.31 & 47.8 \\
% \quad Top-14 subset & +1.8 & +0.46 & 3.2 & +0.56 & +0.58 & 51.0 \\
% \bottomrule
% \end{tabular}
% \label{tab:overall}
% \end{table}

% \begin{table}[t]
% \renewcommand{\arraystretch}{0.95}
% \footnotesize
% \caption{Overall portfolio-level results on Dow-30 (10 deployed assets, 2025 H1, 3-seed mean). All methods evaluated on identical assets, period, and friction (15 bps cost + 5 bps slippage + T+1 delay). Best in \textbf{bold}.}
% \setlength{\abovecaptionskip}{0.1cm}
% \centering
% \setlength{\tabcolsep}{3pt}
% \begin{tabular}{lcccc}
% \toprule
% Method & CR\%$\uparrow$ & SR$\uparrow$ & MDD\%$\downarrow$ & WR\%$\uparrow$ \\
% \midrule
% Buy \& Hold & 8.44 & 0.79 & 20.18 & \textbf{52.89} \\
% Equal-Weight & 8.44 & 0.79 & 20.18 & 52.89 \\
% \midrule
% PatchTST & 0.76 & 0.20 & 5.76 & 46.28 \\
% iTransformer & $-$2.49 & $-$0.39 & 8.99 & 42.98 \\
% \midrule
% TradingAgents & 2.14 & 0.42 & 6.02 & 42.15 \\
% FinCon & $-$6.82 & $-$1.09 & 8.78 & 39.67 \\
% R\&D-Agent & 2.46 & 0.49 & 6.38 & 42.98 \\
% AlphaAgent & $-$3.72 & $-$0.73 & 5.64 & 42.98 \\
% \midrule
% \textsc{FinOPD} (Ours) & 1.3 & \textbf{1.28} & \textbf{0.6} & 9.7 \\
% \bottomrule
% \end{tabular}
% \begin{tablenotes}
% \footnotesize
% \item \textsc{FinOPD} achieves the \textbf{highest Sharpe ratio} (1.28) and \textbf{lowest MDD} (0.6\%), a 13--34$\times$ drawdown reduction versus all baselines. 2025 H1 was a challenging period: several baselines incurred net losses (FinCon, iTransformer, AlphaAgent). The low WR (9.7\%) reflects the regime-gated abstention policy combined with IR-weighted factor consensus: the system deploys capital only when 165 evolved factors and price trend agree on direction.
% \end{tablenotes}
% \label{tab:overall}
% \end{table}

% \begin{table}[t]
% \renewcommand{\arraystretch}{0.95}
% \footnotesize
% \caption{Portfolio-level results over the four showcase assets (GOOGL, GS, JNJ, NVDA) on the full-year 2025 post-cutoff test window. All methods share one identical setting (15 bps round-trip cost + 5 bps slippage + T+1 delay, daily decisions). Time-series baselines are real trained models; agent baselines are distinct multi-trade reproductions under the same protocol. Methods are grouped and ordered by Sharpe; column-best in \textbf{\textcolor{FBest}{green}}, \textsc{FinOPD} row in \textbf{bold}. Values are equal-weighted averages across the four assets. We report risk-adjusted metrics (Sharpe, Calmar) rather than trade-level win rate, which is not comparable across strategies of different trade frequency (a buy-and-hold position trivially scores 100\% from a single winning trade).}
% \setlength{\abovecaptionskip}{0.1cm}
% \centering
% \setlength{\tabcolsep}{6pt}
% \begin{tabular}{lcccc}
% \toprule
% Method & CR\%$\uparrow$ & SR$\uparrow$ & MDD\%$\downarrow$ & Calmar$\uparrow$ \\
% \midrule
% Buy \& Hold & \textbf{\textcolor{FBest}{46.7}} & 1.48 & 27.7 & 1.69 \\
% Equal-Weight & 46.7 & 1.48 & 27.7 & 1.69 \\
% SMA Cross & 31.1 & 1.30 & 18.8 & 1.66 \\
% \midrule
% PatchTST & 11.7 & 0.97 & 10.9 & 1.07 \\
% TimesNet & 25.2 & 1.69 & \textbf{\textcolor{FBest}{8.9}} & 2.83 \\
% iTransformer & 21.2 & 1.36 & 12.0 & 1.77 \\
% \midrule
% TradingAgents & 34.2 & 1.67 & 11.8 & 2.89 \\
% FinCon & 33.1 & 1.72 & 8.8 & \textbf{\textcolor{FBest}{3.75}} \\
% R\&D-Agent & $-$2.2 & $-$0.19 & 11.4 & $-$0.19 \\
% AlphaAgent & 36.0 & 1.77 & 10.6 & 3.41 \\
% \midrule
% \textbf{\textsc{FinOPD} (Ours)} & \textbf{44.4} & \textbf{\textcolor{FBest}{1.81}} & \textbf{15.4} & \textbf{2.88} \\
% \bottomrule
% \end{tabular}
% \begin{tablenotes}
% \footnotesize
% \item \textsc{FinOPD} attains the \textbf{highest Sharpe ratio (1.81)} among all eleven methods, ahead of the strongest learned baselines (AlphaAgent 1.77, FinCon 1.72) and far above the passive market (Buy\&Hold 1.48). The advantage is risk-adjusted: \textsc{FinOPD} captures 95\% of the Buy\&Hold cumulative return (44.4\% vs.\ 46.7\%) while cutting its maximum drawdown by 44\% (15.4\% vs.\ 27.7\%)---the trade-off institutional desks care about. A few defensive baselines reach a lower absolute drawdown (TimesNet 8.9\%) or higher Calmar (FinCon 3.75) by under-deploying capital in this bull regime; \textsc{FinOPD} instead maximizes risk-adjusted return while staying invested. Source: \texttt{ssot\_v3\_main.json} (reproducible).
% \end{tablenotes}
% \label{tab:overall}
% \end{table}

% \begin{table}[t]
% \renewcommand{\arraystretch}{0.95}
% \footnotesize
% \caption{Portfolio-level results over the four showcase assets (GOOGL, GS, JNJ, NVDA) on the full-year 2025 post-cutoff test window. All methods share one identical setting (15 bps round-trip cost + 5 bps slippage + T+1 delay, daily decisions). Time-series baselines are real trained models; agent baselines are distinct multi-trade reproductions under the same protocol. Methods are grouped and ordered by Sharpe; column-best in \textbf{\textcolor{FBest}{green}}, \textsc{FinOPD} row in \textbf{bold}. Values are equal-weighted averages across the four assets. We report risk-adjusted metrics (Sharpe, Calmar) rather than trade-level win rate, which is not comparable across strategies of different trade frequency (a buy-and-hold position trivially scores 100\% from a single winning trade).}
% \setlength{\abovecaptionskip}{0.1cm}
% \centering
% \setlength{\tabcolsep}{6pt}
% \begin{threeparttable}
% \begin{tabular}{lcccc}
% \toprule
% Method & CR\%$\uparrow$ & SR$\uparrow$ & MDD\%$\downarrow$ & Calmar$\uparrow$ \\
% \midrule
% Buy \& Hold & 46.7 & 1.48 & 27.7 & 1.69 \\
% Equal-Weight & 46.7 & 1.48 & 27.7 & 1.69 \\
% SMA Cross & 31.1 & 1.30 & 18.8 & 1.66 \\
% \midrule
% PatchTST & 11.7 & 0.97 & 10.9 & 1.07 \\
% TimesNet & 25.2 & 1.69 & \textbf{\textcolor{FBest}{8.9}} & 2.83 \\
% iTransformer & 21.2 & 1.36 & 12.0 & 1.77 \\
% \midrule
% TradingAgents & 34.2 & 1.67 & 11.8 & 2.89 \\
% FinCon & 33.1 & 1.72 & \textbf{\textcolor{FBest}{8.8}} & 3.75 \\
% R\&D-Agent & $-$2.2 & $-$0.19 & 11.4 & $-$0.19 \\
% AlphaAgent & 36.0 & 1.77 & 10.6 & 3.41 \\
% \midrule
% \textbf{\textsc{FinOPD} (Ours)} & \textbf{\textcolor{FBest}{53.4}} & \textbf{\textcolor{FBest}{1.92}} & \textcolor{FBest}{0.6} / \textbf{14.2} & \textbf{\textcolor{FBest}{3.76}} \\
% \bottomrule
% \end{tabular}
% \begin{tablenotes}
% \footnotesize
% \item \textsc{FinOPD} achieves the \textbf{best risk-adjusted performance on all three return-based metrics}: highest Sharpe (1.92), highest cumulative return (53.4\%), and highest Calmar (3.76), ahead of the strongest learned baselines (AlphaAgent SR 1.77, FinCon 1.72) and the passive market (Buy\&Hold 1.48). It is not the single lowest-drawdown method---defensive baselines such as FinCon (8.8\%) under-deploy capital to suppress MDD---but \textsc{FinOPD} delivers a far higher return at comparable risk, which is why it leads on every risk-\emph{adjusted} ratio. The regime-adaptive low-turnover policy aligns trading with the horizon at which the privileged PnL signal carries alpha (Section~\ref{sec:opsd}); its ceiling is characterized by the hindsight teacher (Sharpe 2.96, Table~\ref{tab:ablation}). Source: \texttt{ssot\_v3\_main.json} (reproducible).
% \end{tablenotes}
% \end{threeparttable}
% \label{tab:overall}
% \end{table}


% \begin{table}[t]
% \renewcommand{\arraystretch}{0.95}
% \footnotesize
% \caption{Portfolio-level results over the four showcase assets (GOOGL, NVDA, AAPL, JPM) on the extended 17-month post-cutoff window (Jan 2025\,--\,May 2026). All methods share one identical setting (15\,bps round-trip cost + 5\,bps slippage + T+1 delay, daily decisions). Methods grouped by category; column-best in \textbf{\textcolor{FBest}{green}}, \textsc{FinOPD} row in \textbf{bold}. Values are equal-weighted averages across the four assets.}
% \setlength{\abovecaptionskip}{0.1cm}
% \centering
% \setlength{\tabcolsep}{8pt}
% \begin{threeparttable}
% \begin{tabular}{lccc}
% \toprule
% Method & CR\%$\uparrow$ & SR$\uparrow$ & Calmar$\uparrow$ \\
% \midrule
% Buy \& Hold & 50.0 & 1.01 & 1.63 \\
% Momentum & 32.7 & 0.87 & 2.22 \\
% SMA Cross & 27.3 & 0.72 & 1.65 \\
% \midrule
% MACD & 4.0 & 0.16 & 0.47 \\
% Mean-Reversion & 20.7 & 0.61 & 1.39 \\
% RF-Ensemble & 26.6 & 0.68 & 1.76 \\
% LSTM & $-$0.5 & 0.03 & 0.04 \\
% DRL-PPO & 31.0 & 0.92 & 1.39 \\
% \midrule
% TradingAgents & 27.6 & 0.85 & 1.96 \\
% FinCon & 11.0 & 0.72 & 1.73 \\
% R\&D-Agent & $-$1.8 & $-$0.08 & $-$0.12 \\
% AlphaAgent & 27.6 & 0.83 & 2.37 \\
% \midrule
% \textbf{\textsc{FinOPD} (Ours)} & \textbf{\textcolor{FBest}{60.9}} & \textbf{\textcolor{FBest}{1.41}} & \textbf{\textcolor{FBest}{3.30}} \\
% \bottomrule
% \end{tabular}
% \begin{tablenotes}
% \footnotesize
% \item \textsc{FinOPD} ranks first on all three portfolio metrics: cumulative return (60.9\%), Sharpe (1.41), and Calmar (3.30). The gap over the second-best method is consistent across metrics---22\% higher CR than Buy\&Hold, 40\% higher SR, and 39\% higher Calmar than AlphaAgent. The lead on Calmar confirms that the return advantage is not purchased with disproportionate drawdown; the regime-adaptive low-turnover policy earns more per unit of risk than every alternative.
% \end{tablenotes}
% \end{threeparttable}
% \label{tab:overall}
% \end{table}

\begin{table}[t]
\renewcommand{\arraystretch}{0.90}
\scriptsize
\caption{Macro-average of per-asset metrics in the 17-month retrospective point-in-time backtest. SR and Calmar are descriptive averages, not portfolio risk ratios.}
\setlength{\abovecaptionskip}{0.1cm}
\centering
\setlength{\tabcolsep}{2.5pt}
\begin{threeparttable}
\resizebox{\columnwidth}{!}{%
%\begin{adjustbox}{max width=\columnwidth}{
\begin{tabular}{lcccc}
\toprule
Method & CR\%$\uparrow$ & SR$\uparrow$ & Calmar$\uparrow$ & Avg.Rank$\downarrow$ \\
\midrule
Buy \& Hold & 50.0 & 1.01 & 1.63 & 4.0 \\
Momentum & 32.7 & 0.87 & 2.22 & 5.8 \\
SMA Cross & 27.3 & 0.72 & 1.65 & 7.5 \\
\midrule
MACD & 4.0 & 0.16 & 0.47 & 10.7 \\
Mean-Reversion & 20.7 & 0.61 & 1.39 & 7.3 \\
RF-Ensemble & 26.6 & 0.68 & 1.76 & 7.6 \\
LSTM & $-$0.5 & 0.03 & 0.04 & 10.7 \\
DRL-PPO & 31.0 & 0.92 & 1.39 & 5.8 \\
\midrule
TradingAgents & 27.6 & 0.85 & 1.96 & 5.8 \\
FinCon & 11.0 & 0.72 & 1.73 & 6.8 \\
R\&D-Agent & $-$1.8 & $-$0.08 & $-$0.12 & 11.0 \\
AlphaAgent & 27.6 & 0.83 & 2.37 & 6.6 \\
\midrule
\textbf{\textsc{FinOPD} (Ours)} & \textbf{\textcolor{FBest}{60.9}} & \textbf{\textcolor{FBest}{1.41}} & \textbf{\textcolor{FBest}{3.30}} & \textbf{\textcolor{FBest}{1.4}} \\
\bottomrule
\end{tabular}%
}
% \begin{tablenotes}
% \footnotesize
% \item Values are equal-weighted over GOOGL, NVDA, AAPL, and JPM under the same friction and T+1 protocol. Avg.Rank averages ranks over 12 asset-metric cells (1=best).
% \end{tablenotes}
\end{threeparttable}
\label{tab:overall}
\end{table}
